
%==========================
% mandopop_proposal.tex
%==========================
\documentclass[11pt]{article}

\usepackage[margin=1in]{geometry}
\usepackage{setspace}
\setstretch{1.0}

\usepackage[T1]{fontenc}
\usepackage{newtxtext,newtxmath}
\usepackage{microtype}
\usepackage[hidelinks]{hyperref}
\usepackage[round,authoryear]{natbib}
\usepackage{enumitem}
\usepackage{titlesec}

\setlength{\parindent}{0pt}
\setlength{\parskip}{0.45em}
\titlespacing*{\section}{0pt}{0.7\baselineskip}{0.25\baselineskip}

\newcommand{\CandidateName}{Frank Yin} % <-- edit
\newcommand{\DegreeProgram}{PhD Proposal (Music Theory)} % <-- edit
\newcommand{\ProposalDate}{\today} % <-- edit

\begin{document}

\begin{center}
{\Large \textbf{Platformized Form and Musical Homogenization in Mandopop (2000--2020):}\\
\textbf{An Audio-Based Corpus Study with Interpretive Media/Gender Analysis}}\\[0.6em]
\CandidateName 

% \quad $\cdot$ \quad \DegreeProgram \quad $\cdot$ \quad \ProposalDate
\end{center}

\vspace{-0.25em}
\section*{1.\ Introduction}
This project proposes an empirical, audio-based corpus study of Mandopop between 2000 and 2020, with a music-theoretical emphasis on \emph{form} (harmony, sectional design, phrase periodicity, repetition strategies) and complementary measurements of melody (motive-wise), and timbre/production, with the core goal to operationalize and test claims about increasing \emph{homogeneity} (or decreasing diversity) in measurable musical features over time.

The historical frame begins near China’s WTO accession and ends at 2020 to avoid confounding the analysis with post-pandemic shocks. I assert 2009, 2012, and 2018 as \emph{candidate} inflection points motivated by changes in China’s platform/media ecology and governance narratives; asserting them as hypotheses, I will test them using change-point methods and compare them to data-driven breakpoints \citep{mauch2015,wang2021platformization,li2025p2p,creemers2020sovereignty}.

Methodologically, the project synthesizes: (i) music-theoretical corpus traditions in popular-music research \citep{leveillegauvin2015,duinker2020,shea2020diss}; (ii) state-of-the-art but \emph{off-the-shelf} music information retrieval (MIR) tools for melody extraction, beat/downbeat tracking, and structural segmentation \citep{salamon2012melody,bock2016madmom,mcfee2015librosa,mcfee2014structure,nieto2020structure}; and (iii) explicit bias auditing and sampling transparency in corpus formation \citep{shea2022sampling,shea2024diversity}. An interpretive layer asks how measured musical shifts can be read through debates about platformization, digital sovereignty, and gender politics---including whether changes correlate with constraint/oppression or new affordances for women in Mandopop’s star/producer ecology \citep{fungcurtin2002faye,moskowitz2010,dekloet2010}.

\section*{2.\ Research questions}
The project answers research questions (RQs), designed to stay anchored in music-analytical evidence while enabling cultural interpretation:

\textbf{RQ1 (Longitudinal change):} How do distributional properties of Mandopop \emph{formal organization} change from 2000--2020 (e.g., section inventories, chorus placement, sectional repetition, phrase periodicity, bridge frequency, sectional duration profiles)? How do these relate to established pop/rock form concepts \citep{covach2005form,nobile2020} and to function-oriented formal thinking (selectively adapted, not transplanted wholesale) \citep{caplin1998}?

\textbf{RQ2 (Homogeneity/diversity):} Do within-year diversity measures and between-year distances show increasing homogeneity in form, melody, and timbre/production? If so, which musical facets drive the trend (e.g., narrowing of sectional prototypes vs.\ narrowing of timbral palettes)? \citep{serra2012,mauch2015}

\textbf{RQ3 (Change points):} Are there statistically robust change points around the hypothesized years 2009, 2012, and 2018? If not, what breakpoints emerge from the musical data, and how stable are they under resampling and alternative corpus constructions? \citep{mauch2015,shea2022sampling}

\textbf{RQ4 (Gendered patterns):} Do formal/timbral trajectories differ by artist gender (and, where possible, by role: vocalist vs.\ producer/composer credits)? Can measured differences be plausibly connected to gendered expectations in Chinese popular music industries, without reducing interpretation to a single causal story? \citep{fungcurtin2002faye,shea2024diversity}

\textbf{RQ5 (Media interpretation):} How can the musical findings be interpreted alongside scholarship on China’s platformization of music circulation and governance discourses around digital/cyber sovereignty? What do ``individualism'' and ``collectivism'' mean as \emph{music-analytical} constructs once operationalized as variation, redundancy, and constraint? \citep{li2025p2p,wang2021platformization,creemers2020sovereignty,dekloet2010}

% comments: not quite audio based, like starting from audio yes, but eventually encoded to symbolic music; shall I address it somewhere?

\section*{3.\ Methodology}
The dissertation applies established music information retrieval (MIR) techniques as analytical instruments rather than as objects of methodological innovation. All analyses are conducted on a stratified corpus of Mandarin-language popular recordings released between 2000 and 2020, sampled evenly across years and major circulation channels to enable longitudinal comparison. Corpus construction, demographic encoding, and representational limitations are documented transparently and evaluated through sensitivity analyses \citep{shea2022sampling,shea2024diversity}.

\textbf{Audio feature extraction.} Audio files are standardized prior to analysis. Melody is estimated using contour-based predominant melody extraction \citep{salamon2012melody}, from which pitch range, interval distributions, contour archetypes, and repetition statistics are derived. Beat and downbeat information is estimated using widely adopted MIR pipelines \citep{bock2016madmom}. Metric interpretations are treated probabilistically, in recognition of analytic disagreement regarding meter in popular music \citep{declercq2016measuring}.

\textbf{Formal analysis.} Musical form is modeled through structural segmentation based on self-similarity and repetition, using spectral clustering approaches \citep{mcfee2014structure,nieto2020structure}. Segment boundaries are converted into theory-facing representations (e.g., verse, chorus, bridge) using repetition strength, timbral stability, and lyric alignment, with limited human correction on a validation subset. Formal interpretation draws selectively on popular-music form scholarship \citep{covach2005form,nobile2020} and on functional concepts adapted from Caplin’s theory of form, understood as interpretive lenses rather than prescriptive templates \citep{caplin1998}.

\textbf{Harmony and timbre.} Harmonic content is represented through chroma-based descriptors and coarse loop typologies, supplemented where reliable by automated chord estimates validated against expert-annotated benchmarks \citep{burgoyne2011chord}. Timbre and production are modeled using spectral features, loudness proxies, and embedding-space representations to capture changes in texture, brightness, and dynamic range.

\textbf{Validation and robustness.} A manually annotated subset (approximately 10\% of the corpus) is used to assess segmentation accuracy, melody extraction reliability, and labeling consistency. All longitudinal claims are accompanied by robustness checks across alternative feature sets and resampling regimes.


\textbf{Analysis and interpretation.}
Homogeneity is quantified through within-year dispersion measures and between-year distributional distances across formal, melodic, harmonic, and timbral features \citep{serra2012,mauch2015}. Change-point analyses compare hypothesized break years (2009, 2012, 2018) with data-driven alternatives. Gender-stratified analyses examine whether musical constraints or affordances evolve differently across performer categories, with interpretation grounded in scholarship on gender politics in Chinese popular music \citep{fungcurtin2002faye,moskowitz2010}.

Rather than asserting direct causality, the dissertation interprets musical change in dialogue with research on platformization and digital sovereignty, asking which musical shifts plausibly align with infrastructural transformation and which resist such narratives \citep{li2025p2p,wang2021platformization,creemers2020sovereignty,dekloet2010}.


\section*{4.\ Contributions}


% \textbf{Contribution:} 

The project contributes (1) the first large-scale, reproducible, audio-based map of Mandopop formal change (2000--2020) centered on form as a primary analytical object; (2) a defensible operationalization of ``homogenization'' that avoids untestable claims about ``quality''; (3) an auditable corpus design and bias report for Mandopop; and (4) a model for integrating MIR-derived evidence with interpretation about platform/media dynamics and gender politics.


\vspace{-0.25em}
\bibliographystyle{plainnat}
\bibliography{mandopop_proposal}

\end{document}



